This report comprehensively reviews the first three chapters of Si Hui Tan's MIT doctoral thesis, ``Quantum State Discrimination with Bosonic Channels and Gaussian States.''\\

We begin by establishing the foundational concepts of quantum state discrimination, including the notion of quantum states, measurement postulates, and the Positive Operator Valued Measure (POVM) formalism. We then explore the two primary discrimination strategies -- minimum-error discrimination (Helstrom bound) and unambiguous discrimination (IDP limit) --along with practical approaches such as the pretty-good measurement and the no-measurement strategy. We highlight the central role of state discrimination in quantum information theory through its applications in quantum key distribution, entanglement characterization, quantum error correction, and quantum algorithms.\\

We further develop the continuous variable quantum information framework, including the quantization of the electromagnetic field, bosonic mode operators, the symplectic structure of phase space, Gaussian transformations, and phase-space representations such as the Wigner function and characteristic function. We provide detailed derivations of coherent states, squeezing operators, two-mode squeezing, and Gaussian EPR states. We introduce Williamson's theorem as a key tool for decomposing covariance matrices and include a complete step-by-step proof.\\

Finally, we examine the distinguishability of Gaussian states using fidelity and the quantum Chernoff bound. We present a numerical implementation using QuTiP to validate the Helstrom bound and the quantum Chernoff bound for multi-copy state discrimination, demonstrating that the quantum Chernoff bound provides an asymptotically tight upper bound on the minimum error probability.
